Equating real parts gives $\cos^2\dfrac{\theta}{2}-\sin^2\dfrac{\theta}{2}=\cos\theta$. Combining this with $\cos^2\dfrac{\theta}{2}+\sin^2\dfrac{\theta}{2}=1$ yields
\[
\cos^2\dfrac{\theta}{2}=\dfrac{1+\cos\theta}{2},\qquad
\sin^2\dfrac{\theta}{2}=\dfrac{1-\cos\theta}{2}.
\]
Thus
\[
\boxed{\cos\dfrac{\theta}{2}=\pm\sqrt{\dfrac{1+\cos\theta}{2}},\quad
\sin\dfrac{\theta}{2}=\pm\sqrt{\dfrac{1-\cos\theta}{2}}}.
\]
For a specified real $\theta$, choose each sign from the quadrant of $\dfrac{\theta}{2}$; the imaginary-part identity $2\sin\dfrac{\theta}{2}\cos\dfrac{\theta}{2}=\sin\theta$ also constrains their product. The two complex square roots are exactly
\[
\boxed{\pm\left(\cos\dfrac{\theta}{2}+i\sin\dfrac{\theta}{2}\right)}.
\]
They differ by a simultaneous sign change, not by arbitrary independent changes to their real and imaginary parts.
\begin{asy}
texpreamble("\usepackage{amsmath,amssymb}");
size(240);
defaultpen(fontsize(9pt));
draw(unitcircle);draw((-1.2,0)--(1.2,0),gray);draw((0,-1.2)--(0,1.2),gray);pair z=dir(70),u=dir(35);draw((0,0)--z,blue);draw(-u--u,red);dot(z,blue);dot(u,red);dot(-u,red);label("$z$",z,N);label("$u$",u,E);label("$-u$",-u,SW);
\end{asy}
